\QHeading{第1题　用简便方法计算下面各题。}
\par\noindent


\CalcBlock[方法一：简便计算]{1}{\begin{aligned}[t]&\phantom{={}}\color{black}23\times5\times2\\&\color{answerred}=23\times(5\times2)\\&\color{answerred}=23\times10\\&\color{answerred}=230\end{aligned}}
\par
\MethodBlock{方法二：按运算顺序直接计算（对照）}{本法只作结果核对；本题要求简便计算，前面的简便解法仍保留。}{\begin{aligned}[t]&\phantom{={}}23\times 5\times 2\\
&=115\times 2\\
&=230\end{aligned}}


\CalcBlock[方法一：简便计算]{2}{\begin{aligned}[t]&\phantom{={}}\color{black}4\times51\times25\\&\color{answerred}=4\times25\times51\\&\color{answerred}=(4\times25)\times51\\&\color{answerred}=100\times51\\&\color{answerred}=5100\end{aligned}}
\MethodBlock{方法二：先处理括号或乘除，再把特殊因数改成整百或整千}{先把25、125等改写成整百或整千除以一个数。}{\begin{aligned}[t]&\phantom{={}}4\times 51\times 25\\
&=204\times 25\\
&=(100\div 4)\times 204\\
&=100\times (204\div 4)\\
&=100\times 51\\
&=5100\end{aligned}}

\CalcBlock[方法一：简便计算]{3}{\begin{aligned}[t]&\phantom{={}}\color{black}125\times24\\&\color{answerred}=125\times(8\times3)\\&\color{answerred}=(125\times8)\times3\\&\color{answerred}=1000\times3\\&\color{answerred}=3000\end{aligned}}
\MethodBlock{方法二：把特殊因数改成整百或整千}{先把25、125等改写成整百或整千除以一个数。}{\begin{aligned}[t]&\phantom{={}}125\times 24\\
&=(1000\div 8)\times 24\\
&=1000\times (24\div 8)\\
&=1000\times 3\\
&=3000\end{aligned}}

\MethodBlock{方法三：按数位拆分计算}{把一个因数拆成整十、整百数与余下部分之和。}{\begin{aligned}[t]&\phantom{={}}125\times 24\\
&=125\times (20+4)\\
&=125\times 20+125\times 4\\
&=2500+500\\
&=3000\end{aligned}}

\CalcBlock[方法一：简便计算]{4}{\begin{aligned}[t]&\phantom{={}}\color{black}2\times13\times50\times2\\&\color{answerred}=2\times50\times13\times2\\&\color{answerred}=(2\times50)\times(13\times2)\\&\color{answerred}=100\times26\\&\color{answerred}=2600\end{aligned}}
\MethodBlock{方法二：结合因数，先凑整}{用乘法交换律、结合律，先算容易凑整的两个因数。}{\begin{aligned}[t]&\phantom{={}}2\times 13\times 50\times 2\\
&=(50\times 2)\times 2\times 13\\
&=100\times 2\times 13\\
&=200\times 13\\
&=2600\end{aligned}}

\MethodBlock{方法三：按运算顺序直接计算（对照）}{本法只作结果核对；本题要求简便计算，前面的简便解法仍保留。}{\begin{aligned}[t]&\phantom{={}}2\times 13\times 50\times 2\\
&=26\times 50\times 2\\
&=1300\times 2\\
&=2600\end{aligned}}

\CalcBlock[方法一：简便计算]{5}{\begin{aligned}[t]&\phantom{={}}\color{black}267+173+127\\&\color{answerred}=267+(173+127)\\&\color{answerred}=267+300\\&\color{answerred}=567\end{aligned}}
\MethodBlock{方法二：按运算顺序直接计算（对照）}{本法只作结果核对；本题要求简便计算，前面的简便解法仍保留。}{\begin{aligned}[t]&\phantom{={}}267+173+127\\
&=440+127\\
&=567\end{aligned}}

\CalcBlock[方法一：简便计算]{6}{\begin{aligned}[t]&\phantom{={}}\color{black}25\times32\\&\color{answerred}=25\times(4\times8)\\&\color{answerred}=(25\times4)\times8\\&\color{answerred}=100\times8\\&\color{answerred}=800\end{aligned}}
\MethodBlock{方法二：把特殊因数改成整百或整千}{先把25、125等改写成整百或整千除以一个数。}{\begin{aligned}[t]&\phantom{={}}25\times 32\\
&=(100\div 4)\times 32\\
&=100\times (32\div 4)\\
&=100\times 8\\
&=800\end{aligned}}

\MethodBlock{方法三：凑整后补偿}{把接近整十、整百的因数凑整，再补回或减去多算的部分。}{\begin{aligned}[t]&\phantom{={}}25\times 32\\
&=25\times (30+2)\\
&=25\times 30+25\times 2\\
&=750+50\\
&=800\end{aligned}}

\MethodBlock{方法四：按数位拆分计算}{把一个因数拆成整十、整百数与余下部分之和。}{\begin{aligned}[t]&\phantom{={}}25\times 32\\
&=32\times (20+5)\\
&=32\times 20+32\times 5\\
&=640+160\\
&=800\end{aligned}}

\CalcBlock[方法一：简便计算]{7}{\begin{aligned}[t]&\phantom{={}}\color{black}540\div45\div2\\&\color{answerred}=540\div(45\times2)\\&\color{answerred}=540\div90\\&\color{answerred}=6\end{aligned}}
\MethodBlock{方法二：先处理括号或乘除，再利用商不变的性质}{被除数和除数同时乘同一个非零数，商不变。}{\begin{aligned}[t]&\phantom{={}}540\div 45\div 2\\
&=12\div 2\\
&=(12\times 5)\div (2\times 5)\\
&=60\div 10\\
&=6\end{aligned}}

\CalcBlock[方法一：简便计算]{8}{\begin{aligned}[t]&\phantom{={}}\color{black}800\div(20\times8)\\&\color{answerred}=800\div20\div8\\&\color{answerred}=40\div8\\&\color{answerred}=5\end{aligned}}
\MethodBlock{方法二：先处理括号或乘除，再拆分除数，连续相除}{}{\begin{aligned}[t]&\phantom{={}}800\div (20\times 8)\\
&=800\div 160\\
&=800\div (2\times 80)\\
&=800\div 2\div 80\\
&=400\div 80\\
&=5\end{aligned}}

\CalcBlock[方法一：简便计算]{9}{\begin{aligned}[t]&\phantom{={}}\color{black}270\div45\\&\color{answerred}=270\div(9\times5)\\&\color{answerred}=270\div9\div5\\&\color{answerred}=30\div5\\&\color{answerred}=6\end{aligned}}
\MethodBlock{方法二：拆分除数，连续相除}{}{\begin{aligned}[t]&\phantom{={}}270\div 45\\
&=270\div (5\times 9)\\
&=270\div 5\div 9\\
&=54\div 9\\
&=6\end{aligned}}

\MethodBlock{方法三：利用商不变的性质}{被除数和除数同时乘同一个非零数，商不变。}{\begin{aligned}[t]&\phantom{={}}270\div 45\\
&=(270\times 2)\div (45\times 2)\\
&=540\div 90\\
&=6\end{aligned}}

\par\noindent

\QHeading{第2题　解决问题。}

\Problem{① 实验小学举行“阳光体育大课间”展示活动，有25个班，每班站4排，每排11人，一共有多少人？}
\par


\MethodTitle{方法一}
\Answer{{\Formula{\begin{aligned}[t]&\phantom{={}}11\times4\times25\\&=11\times(4\times25)\\&=11\times100\\&=1100\text{（人）}\end{aligned}}}\par 答：一共有1100人。}

\MethodTitle{方法二：先求每班人数，再求全校人数}
\Answer{每班4排，每排11人，先求每班44人。}
\Answer{\Formula{\begin{aligned}[t]&\phantom{={}}11\times 4\times 25\\
&=44\times 25\\
&=1100\text{（人）}\end{aligned}}}
\Answer{答：一共有1100人。}
\Problem{② 某食品加工厂上半月用面粉2.36吨，上半月比下半月多用0.25吨，这个食品加工厂全月一共用面粉多少吨？}

\MethodTitle{方法一}
\Answer{{\Formula{\begin{aligned}[t]&\phantom{={}}2.36+(2.36-0.25)\\&=2.36+2.11\\&=4.47\text{（吨）}\end{aligned}}}\par 答：全月用面粉4.47吨。}

\MethodTitle{方法二：先把上、下半月按相同用量计算}
\Answer{上下半月先都按2.36吨算，多算下半月的0.25吨，要减去。}
\Answer{\Formula{\begin{aligned}[t]&\phantom{={}}2.36\times 2-0.25\\
&=4.72-0.25\\
&=4.47\text{（吨）}\end{aligned}}}
\Answer{答：全月一共用面粉4.47吨。}
